% !TeX root = ../papers/nc-p3.tex
\section{平方根の接八次曲面と唯一の長さ二の根}
\label{sec:tangent-determinantal-octics}
有限辺の射影 $X\subset\mathbb P^3$、$D=H^0(\mathcal O_X)$、四文字 $U$、
$Q=U^2$、正規化平方根スキーム $\mathcal F_Q$ を使う。
$U$ の基底を列 $x$ で表す。以下、$X$ が被約な有限辺であるか、
または四辺有限な標準有限Frobenius候補の一位置であるとする。
\begin{lemma}\label{lem:tangent-bordered-octic}
$U$ での接方向は対称四行二次式行列 $H$ であり、$X$ 上で
$H=xh^{\mathsf t}+hx^{\mathsf t}$ を満たすものを取る。
ただし $h$ は $D$ の四成分で、基底変更 $h=Ax$ は除く。
$x$ の張る直線で割った三空間の対称形式として $H$ を読むと、
非零方向の一般階数は2以上である。
さらに $X$ が被約で直線上に四点がなければ一般階数は3であり、
\[
 B=\begin{pmatrix}H&x\\x^{\mathsf t}&0\end{pmatrix},\qquad
 F=-\det B
\]
は $X$ を三重に通る非零八次式である。
\end{lemma}
\begin{proof}
平方の積を一次まで展開すれば表示が得られる。逆に $X$ の各局所環で
$x$ の単元成分を選ぶと、商三空間上で $H=0$ であることから $h$ が一意に解ける。
$H$ が全空間でこの形なら各対角二次式は $x_i$ で割れ、その商の半分は一次式である。
一般点でこの四つの商が $h$ となるので全体でも $h=Ax$ の自明方向に限る。

商上の階数は $(x\wedge -)H(x\wedge -)^{\mathsf t}$ の階数と等しい。
階数1ならこの対称行列を $dgg^{\mathsf t}$ と原始表示する。
$x\wedge g=0$ なので定数二ベクトルの $g$ はあり得ない。$\deg g=1,2$ である。
被約なら $d g_i$ は全て $X$ 上で零となり、次数 $4-\deg g\le3$ の単射性に反する。
非被約ならKoszul完全性より $g=x\wedge v$ で、$v$ は定数または一次の列である。
前者では $d$ は非零二次式で $X$ 上 $d I_{[v]}^2=0$ となり、
定理\ref{thm:nonreduced-opposite-root-separation}の二次式の支持長下界に反する。
後者では定数 $d\ne0$ なので $X$ は $g=x\wedge Ax$ の固有直線スキームの
イデアルの二乗により切られる。非スカラー $A$ でこれは同じ定理のJordan型別長さ評価に反する。

商上で一般階数2なら $B$ は一般階数4で、その随伴の左上四行部分 $P$ は
階数1の非零な対称六次式行列である。$x^{\mathsf t}P=0$ のため
$P=dgg^{\mathsf t}$ で $g$ は定数でない。各成分は $X$ に沿い二重に消える。
単元な $x_i$ を含む二行を消去すると、残りの三行の成分が $I_X$ に入るためである。
被約で直線上に四点がなければ、この六次式行列は存在しない。
定理\ref{thm:opposite-root-triple-octics}の原始次数 $3,2,1$ の各証明は、
二重消滅、三次評価の単射性、直線切断上界だけを使う。
同じ消去で $F$ は $I_X^3$ に入る。
\end{proof}

\begin{theorem}\label{thm:tangent-determinantal-octics}
非零な接方向について $F\ne0$ かつ $B$ の階数3以下の軌跡が有限であるとする。
\begin{enumerate}
\item $U$ での $\mathcal F_Q$ の局所環は $k[t]/(t^2)$ であり、
三次評価同型で全階数の平方根は平方1の単元倍を除いて $U$ だけである。
従って実際の反対座標は全てこれで尽きる。
\item この端を含む第一辺は四辺有限の標準正則延長を高々一つ持ち、
その有無は転置辺との三次交叉と片側対称枝の既出の有限判定で決まる。
反対八次曲面が存在していても成立する。
\item 逆に異なる実際の反対根への四行二次式行列で階数2以下が有限となるなら、
$U$ 自身と相手の根はどちらも正規化平方根スキームの被約孤立点である。
\end{enumerate}
\end{theorem}
\begin{proof}[二十点の二重化の分解]
$S=k[x_0,\ldots,x_3]$ とし、符号をそろえ $P,v$ を $-\operatorname{adj}B$ の
左上四行部分と最後の列の上四成分とする。
$B$ の四次小行列式の分解は対称行列の交代行列、跡零行列、対称行列の複体で得られる。
階数4以上で三写像は期待階数 $10,14,1$ で、階数降下の余次元は3以上なので
定理\ref{thm:double-point-opposite-reconstruction}と同じ完全性判定が使える。
行と列の重みを $(1,1,1,1,0)$ とすれば $I_{2X}$ の極小分解の自由加群のシフトは、
右から次表になる。
\[
\begin{array}{c|c}
\text{生成元}&6^{10},\ 7^4,\ 8^1\\
\text{第一シジジー}&7^4,\ 8^{16},\ 9^4\\
\text{第二シジジー}&9^4,\ 10^6
\end{array}
\]
肩数字は重複度である。重み0同士の成分は零なので定数項はなく極小である。
Hilbert分子は $1-10t^6+15t^8-6t^{10}$ であり、余次元3の飽和イデアルの次数は80である。
$B$ の単元二行の消去により四次小行列式は $I_X^2$ に入り、
局所完全交叉の二重化の長さ80と等しい。よってイデアル層も飽和イデアルも所望のものとなる。

六次生成元は $P$ の十成分で、その一次式係数の関係は $Px=0$ の四本である。
これらだけに関する二次式係数の関係はその一次式倍で尽きる。
実際、次数8の新しい十六関係は
\[
                 HP+xv^{\mathsf t}=F I_4,\qquad x^{\mathsf t}v=F
\]
の右辺 $F$ を消去したものである。任意の定数四行行列 $A$ でこれらを縮約すると、
七次生成元 $v$ の係数は転置規約を除き $(A-\operatorname{tr}(A) I)x$ である。
これが恒等零なら $A=0$ である。したがって七次生成元を使わない新しい次数8の関係はない。
\end{proof}

\begin{proof}[他の根の排除]
全階数平方根の対称行列を $N(x)$ とすると、その六次余因子は $P$ の定数線型結合となる。
十個の二次座標を独立変数に戻し、$Z=(z_{ij})$ を一般対称四行行列、
$J_Z=I_3(Z)$、$J_N=I_3(N(z))$ とする。余因子の対応は十変数環上の次数零の射
\[
                    J_N\longrightarrow J_Z/(\det Z)
\]
を与える。実際 $J_N$ の生成関係は跡零行列による一次関係で尽きる。
それを $x$ 空間に戻すと上の六次生成元だけの二次式関係となり、
$Px=0$ の一次式倍である。十変数に戻したこれらは $P Z=0$ の成分である。
$P$ の代わりに $\operatorname{adj}Z$ を入れると全て $\det Z$ を法として消える。

$N(x)$ の一般階数が3以上なら少なくとも一つの余因子が非零なので射も非零である。
$\det Z=0$ は正規な整超曲面である。特異台は余次元3の階数2以下に限られ、
一般階数3の対称行列の軌道が稠密である。ここでの余次元3は周囲の空間で測る。
この超曲面上で $V(J_N)$ の外では、射は単元イデアルから座標環への射として
正則関数倍となる。除いた部分は超曲面上でも余次元2以上である。
正規性で関数は延長し、次数零であるため定数である。標的が整環のイデアルなので
稠密開での等号は全体に延び、この定数は非零である。
従って $J_N\subset J_Z+(\det Z)=J_Z$、三次生成空間の次元で等号となる。
その特異台のVeroneseが一致するので $N(x)$ は実際には恒等根の表示となり、
一般階数3以上に反する。一般階数2は定理\ref{thm:nonreduced-opposite-root-separation}、
被約有限辺だけの場合は定理\ref{thm:opposite-root-triple-octics}の階数2の証明で除かれる。
一般階数1ならVeroneseが一致し、平方を正規化した単元の相違は符号だけである。
\end{proof}

\begin{proof}[接空間と二次障害]
別の非零接方向 $H'$ に対応する左上余因子 $P'$ は補題の一般階数下界から零でない。
$P'x=0$ は全て $Px=0$ の定数線型結合となる。$x_i$ の一次独立性より係数を比べて
$P'=AP$ を得る。$P$ の十成分は独立で $AP$ は対称なので $A$ はスカラーである。
実際 $A$ とのこの関係を全ての対称行列に対し課せばよい。
$F\ne0$ なので $P$ は一般階数3で、$x$ で割った三次元空間では非零定数因子と
共通の基底の選択を除き $H$ の随伴となる。二つの可逆な三行行列の随伴が
定数比例すれば、元の二行列も定数比例する（比例係数の平方が最初の定数）。
従って $H'$ は同じ接方向の定数倍であり、接空間は一次元である。

二次までの持上げを仮定し、基底を取って平方行列を
$N_t=xx^{\mathsf t}+tH+t^2K$ とする。随伴を
$\operatorname{adj}N_t=t^2 P+t^3 P_1+O(t^4)$ と展開できる。
先頭係数と $\det N_t=t^3 F+O(t^4)$ は、単元な $x_i$ での二次形式の消去、
または階数1行列を加えた行列式の公式で分かる。
二次まで $X$ 上で階数1なので、各 $P_1$ 成分も $I_{2X}(6)$ に入る。
実際単元な対角成分で消去した三行行列は恒等的に定数項が零で、
$t,t^2$ の両係数が $I_X$ に入る。三次小行列式を展開すると
$t^3$ の項はこれらの少なくとも二成分の積からなる。
$N_t\operatorname{adj}N_t=(\det N_t)I$ の $t^3$ 係数を上の境界行列の式と比較すれば
$x^{\mathsf t}P_1=v^{\mathsf t}$ となる。右辺は七次の極小生成元であり、
一次式倍した六次生成元から得られない。これで二次持上げが否定される。
一次元接空間の局所環の完備化は $k[[t]]$ の商であり、二次持上げの不在から長さはちょうど2である。

異なる反対根の四行行列に有限階数降下を仮定した場合は、
定理\ref{thm:double-point-opposite-reconstruction}の分解により $I_{2X}$ の
六次生成元間に一次式係数の関係は一本もない。$U$ に非零接方向があれば
補題により一般階数は2以上で、従ってその境界行列から $P\ne0$ を得る。
$Px=0$ の各行は六次生成元間の一次式関係なので恒等的に係数が零となり、
独立な $x_i$ で係数を比較すれば $P=0$ となる矛盾である。
反対根自身の被約孤立性は定理\ref{thm:nonreduced-opposite-root-separation}の
接方向の証明による。第一辺の延長判定には実際の反対位置が $U$ と射影同値であることから、
系\ref{cor:octic-free-finite-extension}と同じ転置辺の構成を使えばよい。
\end{proof}

\paragraph{正則例における接八次曲面。}
定理\ref{thm:single-edge-instanton-extension}の有理辺で、この条件は実際に成立する。
\path{computations/verify_tangent_determinantal_octics.py} は
$(x\wedge -)H(x\wedge -)^{\mathsf t}$ の四次成分を十五個の小行列式の商へ送り、
良い素数1009で百未知数の核が十七次元であることを検査する。
十六の基底変更と既知の有理シンプレクティック接方向があるので核次元は標数零でも同じで、
この素数で核の基底は持ち上がる。非自明な方向の $F$ は明示平面上で
三偏導の十二次倍が十九次全空間を張り、幾何学的に滑らかな平面切断を持つ。
従って階数3以下は有限であり、この辺の全ての許容反対根の局所スキームも決定される。

\begin{corollary}\label{cor:normal-triple-octic-root-length}
被約有限辺の指定端で直線上に四点がなく、三重消滅八次式の空間が零または
正規な整曲面一本であるとする。全階数かつ三次評価同型の正規化平方根の開集合は、
行の符号倍を除いて、被約一点、長さ二の一点、被約二点のいずれかに尽きる。
七余因子の標識付き反対対 $(X,Y)$ で両向きの標準辺が同じ被約長さ20の
グラフを与え、標準中間座標 $C$ にこの八次・直線条件が成り立つ場合、
全ての四辺有限な標準正則延長の有無は有限に決まる。
延長があればその二つの奇位置は同じ非標準根で、反対階数は9以下である。
\end{corollary}
\begin{proof}
$U$ と異なる根は、被約行列の定理\ref{thm:opposite-root-triple-octics}の証明で
非零行列式を与える。非零接方向も補題\ref{lem:tangent-bordered-octic}で非零八次式を与える。
いずれの表示でも行列の余階数2以上の点は行列式の特異点である。
正規性によりそれらは有限なので、異なる根があれば二重点分解による一意性と
直前の定理の両方の接消滅を適用する。異なる根がなければ接空間は零か、
同じ定理の長さ2の場合になる。符号で再正規化する自由度は三次評価同型での
平方空間の単元安定群がスカラーだけであることによる。

三割平面の問題ではどちらの奇位置も $X$ との完全交叉の跡により $C^2$ の根である。
$C$ 自身はどちらの向きでも定理\ref{thm:trisecant-middle-simple-obstruction}で排除される。
従って両方が唯一の異なる根でなければならない。その座標を $V$ とすれば
反転する二辺 $(V,X,V)$ を持つため対称となり、反対階数は10以下である。
10で正則なら定理\ref{thm:single-edge-instanton-extension}により $V$ での
正規化平方根の接空間に非零方向が必要だが、その根は被約孤立である。
有限判定ではGrassmann上の二次式で全階数・三次評価同型の根を求める。
この開集合は上の記述で有限である。異なる根があれば $(X,V,Y)$ の積の核六式、
四次元三次交叉から候補を作り、四辺有限、有限Frobenius、第四位置の同定を検査する。
階数9以下には系\ref{cor:all-defective-one-sym-finite}を適用すればよい。
\end{proof}
